Para el control de versiones del código fuente y la gestión de la configuración del proyecto se ha utilizado \textbf{Git} como herramienta principal, apoyándose en \textbf{GitHub} como repositorio remoto central. Esta decisión permite mantener un histórico de cambios sobre la aplicación, la documentación, los scripts de apoyo, las migraciones de base de datos y los ficheros de configuración necesarios para reproducir el trabajo.

\subsection{Control de versiones}

Git permite registrar de forma estructurada los cambios realizados durante el desarrollo. Cada modificación queda asociada a un commit, lo que facilita identificar cuándo se introdujo una funcionalidad, una corrección, un ajuste de configuración o una actualización documental. El uso de mensajes descriptivos en los commits ayuda a reconocer el objetivo de cada cambio y a valorar su impacto sobre el resto del sistema.

En este proyecto, al tratarse de un desarrollo individual, no se ha aplicado un flujo complejo de ramas. La rama principal \texttt{master} actúa como referencia integrada del trabajo, tanto para el código fuente como para la memoria académica. Aun así, el uso de Git mantiene la posibilidad de comparar estados, revisar diferencias entre versiones y recuperar un punto anterior si una modificación introduce un error.

\subsection{Gestión de cambios}

Cada cambio relevante se ha tratado como una unidad de trabajo concreta. Antes de considerarlo consolidado, se revisa que la modificación sea coherente con el alcance del proyecto, que no contradiga requisitos o decisiones ya documentadas y que no introduzca dependencias innecesarias respecto a otras partes del sistema.

Esta gestión resulta especialmente útil en un proyecto iterativo, donde algunas decisiones técnicas se han ajustado conforme avanzaba la implementación. El histórico de Git permite conservar la trazabilidad de estas decisiones y diferenciar entre cambios funcionales, correcciones, ajustes de despliegue, modificaciones de documentación y revisiones de configuración.

\subsection{Alcance de la configuración versionada}

El repositorio no controla únicamente el código fuente de la aplicación. También se versionan la documentación en \LaTeX{}, las migraciones de \textit{TypeORM}, los scripts auxiliares y los ficheros de configuración necesarios para construir y ejecutar el sistema. De este modo, se mantiene un conjunto de artefactos suficiente para revisar el proyecto, reproducirlo y continuar su desarrollo sin depender de información externa no documentada.

Quedan fuera del control de versiones las dependencias instaladas, salidas de compilación, cachés, copias temporales, resultados locales de ejecución y ficheros con credenciales o valores sensibles de entorno. Estos elementos se excluyen porque pueden reconstruirse o porque pertenecen a la configuración privada de cada entorno.

\subsection{Migraciones de base de datos}

La evolución del esquema de datos se mantiene mediante migraciones de \textit{TypeORM} \citep{typeormmigrations2026}. Estas migraciones forman parte de la gestión de configuración porque documentan los cambios aplicados sobre la estructura de la base de datos y permiten reconstruirla en un entorno nuevo sin depender de operaciones manuales.

En el estado actual del proyecto existen 48 migraciones versionadas. Su ejecución se realiza mediante el comando:

\begin{lstlisting}[language=bash]
npm run migration:run
\end{lstlisting}

Este mecanismo separa la evolución de la estructura de datos de los datos concretos de cada entorno, manteniendo trazabilidad sobre el modelo persistente y reduciendo el riesgo de inconsistencias entre desarrollo, demostración y futuras instalaciones.

\subsection{Configuración sensible}

Las credenciales, cadenas de conexión, claves privadas y valores específicos de cada entorno no se almacenan directamente en el repositorio. Estos datos se gestionan mediante variables de entorno y ficheros locales excluidos del control de versiones, de forma que la misma base de código pueda ejecutarse en desarrollo, demostración o futuros entornos sin exponer información sensible.

Esta separación entre código y configuración evita publicar secretos junto al proyecto y permite adaptar el sistema a distintos entornos modificando únicamente sus valores de configuración externos. Además, reduce el riesgo de que claves de servicios externos, conexiones a base de datos o parámetros privados queden incorporados accidentalmente al histórico del repositorio.

\subsection{Versiones estables}

Durante el desarrollo académico no se ha establecido una política formal de releases etiquetadas en GitHub. La versión estable de referencia se ha mantenido en la rama \texttt{master}, actualizada de forma incremental conforme se consolidaban funcionalidades, correcciones y cambios documentales.

En una evolución posterior del proyecto, si el sistema pasara a mantenerse como producto reutilizable o desplegado en varios entornos, sería recomendable incorporar etiquetas o releases para identificar versiones cerradas. En el alcance actual del \textit{Trabajo Fin de Grado}, el control mediante commits y rama integrada resulta suficiente para conservar la trazabilidad del desarrollo realizado.
